\subsection{A certified finite development of the Fuchs--Perchik question}
\label{dev:perchik}

\paragraph{The question supported by the notebook.}
On \GSpage{P000655}, Fuchs asks whether nonzero Euler characteristics occur
for arbitrarily large \(m\). The product and the expansion on
\GSpage{P000655}--\GSpage{P000656} make this a definite coefficient question,
not just a request to compute a Lie algebra's cohomology. We write \(E(m)\)
for the coefficients to avoid the notebook's switching between
\(\chi_m\) and \(\chi_m^*\). Put
\[
 S_\ell=\{-\ell-2,-\ell,\ldots,\ell+2\},\qquad
 P(t,x)=\prod_{\ell=1}^{\infty}\prod_{s\in S_\ell}(1-t^\ell x^s),
\]
and define the formal Laurent series
\begin{equation}\label{eq:perchikseries}
 F(t)=\frac12\operatorname{CT}_x
 \bigl[(1-t^{-1}x^{-1})(1-t^{-1}x)
       (1-x^{-2})(1-x^2)P(t,x)\bigr]
 =\sum_{m\geq-1}E(m)t^{2m}.
\end{equation}
Here \(\operatorname{CT}_x\) extracts the coefficient of \(x^0\);
the product is expanded first in nonnegative powers of \(t\) in \(P\),
then multiplied by the displayed finite prefactor. This prescription,
including the two negative-weight factors, is essential.

\begin{question}[The notebook question, precisely normalized]
For every integer \(M\geq0\), is there an integer \(m>M\) with
\(E(m)\ne0\)? Equivalently, is \(F(t)\) not a Laurent polynomial?
\end{question}

\paragraph{The modern target and its limits.}
Let \(\mathfrak h_2\) be the topological complex Lie algebra of formal
Hamiltonian vector fields on the symplectic plane, represented by
\(\mathbb C[[p,q]]/\mathbb C\) with the Poisson bracket. Its quadratic
Hamiltonians form \(\mathfrak{sl}_2\). Use continuous cochains, trivial
coefficients, and input weight \(d-2\) for a Hamiltonian of polynomial
degree \(d\). Then the modern interpretation is
\[
 E(m)=\sum_q(-1)^q
 \dim H^q(\mathfrak h_2,\mathfrak{sl}_2)_{2m}.
\]
The target adopted here is infinite nonvanishing of these \emph{relative}
Euler characteristics at fixed symplectic dimension two. Infinite
nonvanishing of cohomology alone would be weaker: even and odd Betti
numbers might cancel. Passing to a stable symplectic rank, deleting the
constant vector fields, or retaining only positive-degree Hamiltonians
changes the question. No quantitative asymptotic, eventual nonvanishing,
or sign law is asserted.

\paragraph{What the literature establishes.}
The coefficient formula is Perchik's theorem, published in 1976
\citep{perchik1976cohomology}. Harvard's dissertation list dates his
dissertation of the same title to 1975 \citep{perchik1975dissertation}.
That establishes earlier work,
not a precise public-circulation date; the seminar records its discussion
on 12 April 1976. \citet[Theorem~2.1 and the table on Russian pp.~68--69]
{mohammadzadehfuchs2014cohomology} restate the formula, compute through
\(m=129\), and explicitly retain the infinite-nonvanishing question.
Their absolute/relative cohomology comparison excludes weight zero.
As of 10 November 2026, the targeted searches documented with this chapter
located no later resolution of this exact fixed-rank coefficient question.
This is a bounded literature finding, not a proof that no resolution exists.

\paragraph{Why a finite development is useful.}
The notebook's missing star would turn a nontrivial question into an
identically zero absolute Euler-characteristic calculation. Its general
product row also has an uncertain middle exponent. A reproducible expansion
therefore tests a mathematically decisive normalization. The contribution
here is a cutoff proof, two independent exact algorithms, an audited finite
extension of the notebook's table, and explicit bounds on what those data
prove. These are proofs and computations supplied for this companion;
no priority claim is made for the elementary lemmas or numerical values.

\begin{lemma}[Well-definedness, integrality, and parity]\ours\quad
For every integer \(N\geq0\), all coefficients of \(F\) through
\(t^N\) are determined by retaining \(P\) through degree \(N+2\)
and including only factors with \(1\leq\ell\leq N+2\).
The series \(F\) has integer coefficients, no powers below \(-2\),
and no odd powers.
\end{lemma}
\begin{proof}
Each factor in \(P\) has positive \(t\)-degree in its nonconstant term.
Thus a coefficient through degree \(L\) uses only factors with
\(\ell\leq L\); these are finitely many, each with finitely many terms.
The finite prefactor has lowest \(t\)-degree \(-2\), so \(L=N+2\)
suffices. Terms discarded beyond \(L\) cannot return to a degree at most
\(N\).

Set \(R=(1-t^{-1}x^{-1})(1-t^{-1}x)P\). Each coefficient of \(R\)
is an integer Laurent polynomial invariant under \(x\mapsto x^{-1}\).
Writing such a coefficient as \(Q(x)=\sum_k a_kx^k\), with
\(a_{-k}=a_k\), gives
\[
 \tfrac12\operatorname{CT}_x[(1-x^{-2})(1-x^2)Q]
   =\tfrac12(2a_0-a_2-a_{-2})=a_0-a_2\in\mathbb Z.
\]
Every monomial \(t^n x^k\) in \(P\) satisfies \(n\equiv k\pmod2\),
because \(s\equiv\ell\pmod2\) in each factor. The negative-weight
factors and the two Weyl factors preserve this congruence. Taking the
constant term therefore kills every odd \(t\)-degree. Finally, the only
negative \(t\)-degrees come from the two factors of degree \(-1\).
\end{proof}

\begin{remark}[Why this is the relative Euler series]\ours\quad
This interpretation can be checked independently at the level of characters.
As an \(\mathfrak{sl}_2\)-module, the weight-\(\ell\) Hamiltonians are
\(\operatorname{Sym}^{\ell+2}\mathbb C^2\), whose weights are \(S_\ell\).
The relative exterior complex removes the entire quadratic component and
retains \(\operatorname{Sym}^{1}\mathbb C^2\) in weight \(-1\) and
the components in positive weights. Its alternating character is \(R\).
For an irreducible character
\(x^d+x^{d-2}+\cdots+x^{-d}\), the difference between its coefficients
of \(x^0\) and \(x^2\) is one for \(d=0\) and zero otherwise.
Thus the operation in the proof selects the invariant multiplicity.
The chain groups at a fixed input weight are finite: at most two
negative-weight generators can occur, and the number and weights of
positive-weight generators are bounded by the total weight plus two.
Euler--Poincare therefore applies weight by weight. This derivation also
explains why an unrestricted exterior product containing the zero-weight
Cartan direction would give the factor \(1-1\), not the desired series.
\end{remark}

\begin{proposition}[Two exact coefficient algorithms]\ours\quad
Write \(P(t,x)=\sum_{n\geq0}P_n(x)t^n\), with \(P_0=1\), and put
\[
 B_r(x)=-\sum_{d\mid r}d\sum_{s\in S_d}x^{sr/d}.
\]
Then, for \(n\geq1\),
\begin{equation}\label{eq:perchikrecurrence}
 nP_n(x)=\sum_{r=1}^n B_r(x)P_{n-r}(x).
\end{equation}
Independently, multiplying each factor \(1-t^\ell x^s\), with descending
updates in \(n\), computes the same coefficients using integer additions
and subtractions only. Both methods terminate at any prescribed cutoff.
\end{proposition}
\begin{proof}
The positive product has constant term one. Its formal logarithmic
derivative in \(\mathbb Q[x,x^{-1}][[t]]\) is
\[
 \frac{t\partial_t P}{P}
 =-\sum_{\ell\geq1}\sum_{s\in S_\ell}
       \frac{\ell t^\ell x^s}{1-t^\ell x^s}
 =-\sum_{\ell\geq1}\sum_{s\in S_\ell}\sum_{j\geq1}
       \ell t^{\ell j}x^{sj}
 =\sum_{r\geq1}B_r(x)t^r.
\]
All coefficient rearrangements are finite at a fixed \(t\)-degree.
Multiplying by \(P\) and comparing the coefficient of \(t^n\) proves
\eqref{eq:perchikrecurrence}; division by \(n\) determines \(P_n\)
uniquely in characteristic zero. For the other algorithm, if \(D_{n,k}\)
is the coefficient before one factor is applied, the new coefficient is
\(D_{n,k}-D_{n-\ell,k-s}\). Descending \(n\) ensures that the second
term still belongs to the old product. Induction over factors proves the
algorithm. The cutoff lemma proves termination and prevents omitted
factors from changing any requested coefficient. Integrality of all
divisions in the recurrence follows also from its agreement with the
integer product; the implementation checks every remainder directly.
\end{proof}

\paragraph{Independent checks of the argument.}
The cutoff was re-derived from exterior degrees: the negative component
has dimension two, hence can lower a positive weight by at most two.
The invariant-multiplicity argument above independently checks the
factor \(1/2\) and both Weyl factors. The recurrence was checked by
direct factor multiplication, rather than by another implementation of
the same recurrence. In particular, agreement is checked for every
Laurent coefficient of every \(P_n\), not just for their final constant
terms. These checks found no mathematical discrepancy between the two
algorithms.

\begin{proposition}[Finite output through \(m=150\)]\oursComputed\quad
The series in \eqref{eq:perchikseries}, computed with exact integers by
both algorithms through positive-product degree \(302\), has
\(135\) nonzero coefficients among \(E(1),\ldots,E(150)\), and
\[
 \sum_{m=1}^{150}|E(m)|=9\,701\,705\,312.
\]
In particular,
\[
 E(35)=10,\quad E(86)=-22\,044,\quad
 E(129)=7\,151\,485,\quad E(150)=1\,408\,548\,898.
\]
\end{proposition}
\begin{proof}
The executable certificate is \texttt{code/perchik\_series.py}; the complete
coefficient table and verification output are in
\texttt{computations/m150/}. The two algorithms agree for all Laurent
coefficients through \(n=302\), including \(137\,561\) checked integer
divisions in the recurrence. Reflection, parity, and the support bound
\(|k|\leq3n\) are checked for every nonzero coefficient. Separate runs
at cutoffs \(82\) and \(182\) produce identical overlapping output.
The notebook's displayed terms, including \(10t^{70}\), are recovered.
The companion verification script records the prefix comparisons and
the sums in the statement.
\end{proof}

\begin{center}
\begin{tabular}{rr@{\qquad}rr@{\qquad}rr}
\toprule
\(m\)&\(E(m)\)&\(m\)&\(E(m)\)&\(m\)&\(E(m)\)\\
\midrule
130&41831534&137&-67794447&144&-30506945\\
131&79675254&138&-188663902&145&321390415\\
132&113436381&139&-311827536&146&732255863\\
133&133851464&140&-412439649&147&1137936570\\
134&131414489&141&-459062275&148&1450364576\\
135&98557132&142&-423731634&149&1571713579\\
136&31149174&143&-282894462&150&1408548898\\
\bottomrule
\end{tabular}
\end{center}
This table is generated by the supplied program; it is not a transcription
of a literature table. It extends the notebook's range, not a claim about
the largest range ever computed.

\begin{corollary}[Exactly what nonzero Euler characteristics prove]\ours\quad
For the relative continuous cohomology above,
\[
 \sum_{m=1}^{150}\sum_q
    \dim H^q(\mathfrak h_2,\mathfrak{sl}_2)_{2m}
       \geq9\,701\,705\,312.
\]
At least \(135\) distinct positive weights in this range carry nonzero
relative cohomology. Also
\(\sum_{m=-1}^{35}|E(m)|=57\), recovering the notebook's finite
lower-bound count with the relative interpretation.
\end{corollary}
\begin{proof}
For each finite weight complex, Euler--Poincare and the triangle
inequality give \(|E(m)|\leq\sum_q\dim H^q_{2m}\). Sum this inequality
over the stated ranges. Distinct weights give distinct graded summands.
No inference is made at a weight whose Euler characteristic vanishes.
\end{proof}

\paragraph{The remaining obstruction and the next bounded step.}
The equality of two exact algorithms certifies a finite prefix. It cannot
exclude a later vanishing tail, prove infinitely many sign changes,
recover the degrees of the cohomology classes, or identify representatives.
The obstacle is cancellation in a two-variable exterior product after
projection to invariants; the large unsigned chain dimensions do not
provide a lower bound on its signed constant term. The checked literature
shows continuing computation rather than abandonment; it does not
establish a sociological explanation for the absence of a general proof.
A further bounded run could compute a longer independently certified
prefix or construct and reduce the relative differential at one selected
weight with \(E(m)=0\), thereby detecting cancellations invisible to Euler
characteristics. A resolution of GP-1975-20 needs an argument excluding
eventual cancellation for unbounded \(m\); extending the table alone
does not supply that argument.
